\frontmatter

% ============================================================
% Halaman Sampul
% ============================================================
\begin{titlepage}
  \centering
  \vspace*{1cm}
  {\LARGE\bfseries Pemrograman Game\par}
  \vspace{0.5cm}
  {\large Menggunakan Unity Engine dan C\#\par}
  \vspace{0.5cm}
  {\large Buku Ajar Berbasis Outcome-Based Education (OBE)\par}
  \vspace{1cm}
  {\large Mata Kuliah: Pemrograman Game\par}
  {\large Kode: IF-XXX\par}
  \vspace{2cm}
  {\large Program Studi S1 Informatika\par}
  {\large Fakultas Teknologi Informasi\par}
  {\large Universitas Lorem Ipsum\par}
  \vfill
  {\large 2026\par}
\end{titlepage}

\cleardoublepage

% ============================================================
% Kata Pengantar
% ============================================================
\chapter*{Kata Pengantar}
\addcontentsline{toc}{chapter}{Kata Pengantar}

Puji syukur kehadirat Tuhan Yang Maha Esa atas terselesaikannya buku ajar \textit{Pemrograman Game} ini. Buku ini disusun dengan pendekatan \textbf{Outcome-Based Education (OBE)}, yang berfokus pada pencapaian kompetensi terukur sesuai dengan Capaian Pembelajaran Lulusan (CPL) dan Capaian Pembelajaran Mata Kuliah (CPMK).

Pemrograman Game merupakan bidang yang sangat menarik dan relevan dalam era digital saat ini. Unity Engine sebagai salah satu game engine terpopuler menyediakan platform yang powerful untuk mengembangkan berbagai jenis game, dari game sederhana hingga game AAA. Dengan kombinasi Unity dan C\#, mahasiswa dapat mempelajari konsep-konsep fundamental game development sekaligus mengasah keterampilan programming.

Buku ini dirancang untuk mendukung pembelajaran mahasiswa dengan pendekatan student-centered learning. Setiap bab dilengkapi dengan:
\begin{itemize}
  \item Sub-CPMK yang jelas dan terukur
  \item Materi pokok dengan contoh praktis di Unity
  \item Aktivitas pembelajaran yang mendorong eksplorasi mandiri
  \item Latihan dan refleksi untuk penguatan pemahaman
  \item Asesmen untuk mengukur pencapaian kompetensi
  \item Checklist kompetensi untuk self-assessment
\end{itemize}

Kami berharap buku ini dapat menjadi panduan yang efektif dalam perjalanan pembelajaran Anda menguasai Pemrograman Game dengan Unity Engine.

\vspace{1cm}
\begin{flushright}
Penyusun\\
Dr. Lorem Ipsum, M.Kom.
\end{flushright}

\cleardoublepage

% ============================================================
% Cara Menggunakan Buku Ini
% ============================================================
\chapter*{Cara Menggunakan Buku Ini}
\addcontentsline{toc}{chapter}{Cara Menggunakan Buku Ini}

Buku ajar ini dirancang dengan pendekatan OBE untuk memaksimalkan pencapaian pembelajaran Anda. Berikut panduan penggunaan buku ini:

\section*{Struktur Buku}

\textbf{Bab I: Pendahuluan dan Orientasi}\\
Memperkenalkan tujuan buku, keterkaitan dengan RPS, dan konteks kurikulum OBE.

\textbf{Bab II: Landasan Teori}\\
Menyajikan fondasi teoretis game development yang menjadi basis pembelajaran seluruh bab berikutnya.

\textbf{Bab III-XV: Unit Materi Inti}\\
Setiap bab mencakup satu topik utama game development dengan struktur lengkap: Sub-CPMK, materi, aktivitas, latihan, asesmen, dan checklist.

\textbf{Bab XVI: Evaluasi dan Refleksi Akhir}\\
Berisi asesmen komprehensif dan panduan refleksi untuk mengukur pencapaian kompetensi secara menyeluruh.

\textbf{Lampiran}\\
Menyediakan rubrik penilaian, glosarium istilah game development, dan referensi tambahan.

\section*{Komponen dalam Setiap Bab}

\begin{enumerate}
  \item \textbf{Sub-CPMK}: Baca dengan seksama untuk memahami kompetensi yang harus dicapai
  \item \textbf{Materi Pokok}: Pelajari dengan cermat, praktikkan semua contoh di Unity
  \item \textbf{Aktivitas Pembelajaran}: Lakukan secara mandiri atau berkelompok
  \item \textbf{Latihan}: Kerjakan untuk menguji pemahaman Anda
  \item \textbf{Asesmen}: Gunakan untuk mengukur pencapaian Sub-CPMK
  \item \textbf{Checklist}: Centang setelah yakin menguasai setiap indikator
\end{enumerate}

\section*{Tips Belajar Efektif}

\begin{itemize}
  \item Jangan hanya membaca, praktikkan semua contoh di Unity Editor
  \item Gunakan Unity Hub untuk mengelola versi dan proyek
  \item Kerjakan latihan sebelum melihat solusi
  \item Diskusikan konsep yang sulit dengan teman atau dosen
  \item Manfaatkan checklist untuk self-assessment berkala
  \item Kerjakan proyek mini untuk mengintegrasikan konsep yang dipelajari
\end{itemize}

\cleardoublepage

% ============================================================
% Identitas Mata Kuliah
% ============================================================
\chapter*{Identitas Mata Kuliah}
\addcontentsline{toc}{chapter}{Identitas Mata Kuliah}

\begin{tabular}{ll}
  Nama Program Studi & : S1 Informatika \\
  Nama Mata Kuliah & : Pemrograman Game \\
  Kode Mata Kuliah & : IF-XXX \\
  Semester & : 5 (Lima) \\
  SKS / Bobot Kredit & : 3 SKS (2 Teori, 1 Praktikum) \\
  Dosen Pengampu & : Dr. Lorem Ipsum, M.Kom. \\
  Tanggal Penyusunan & : 7 Februari 2026 \\
\end{tabular}

\vspace{1cm}

\section*{Capaian Pembelajaran Lulusan (CPL)}

CPL yang dibebankan pada mata kuliah ini mencakup kompetensi lulusan dalam aspek pengetahuan, keterampilan, dan sikap:

\begin{enumerate}
  \item \textbf{CPL-1 (Pengetahuan):} Menguasai konsep teoretis bidang pengetahuan tertentu secara umum dan konsep teoretis bagian khusus dalam bidang pengetahuan tersebut secara mendalam, serta mampu memformulasikan penyelesaian masalah prosedural.
  
  \item \textbf{CPL-2 (Keterampilan Umum):} Mampu menerapkan pemikiran logis, kritis, sistematis, dan inovatif dalam konteks pengembangan atau implementasi ilmu pengetahuan dan teknologi yang memperhatikan dan menerapkan nilai humaniora.
  
  \item \textbf{CPL-3 (Keterampilan Khusus):} Mampu merancang, mengimplementasikan, dan mengevaluasi game menggunakan Unity Engine dan C\# dengan mempertimbangkan standar kualitas software game.
  
  \item \textbf{CPL-4 (Sikap):} Menunjukkan sikap bertanggung jawab atas pekerjaan di bidang keahliannya secara mandiri dan mampu bekerja sama dalam tim.
\end{enumerate}

\section*{Capaian Pembelajaran Mata Kuliah (CPMK)}

Kemampuan atau kompetensi spesifik yang diharapkan mahasiswa kuasai setelah menyelesaikan mata kuliah:

\begin{enumerate}
  \item \textbf{CPMK-1:} Mahasiswa mampu memahami dan menjelaskan konsep dasar game development (game loop, game engine, GameObject, component).
  
  \item \textbf{CPMK-2:} Mahasiswa mampu menggunakan Unity Engine untuk membuat scene, mengelola GameObject, dan mengimplementasikan basic gameplay mechanics.
  
  \item \textbf{CPMK-3:} Mahasiswa mampu mengimplementasikan game logic menggunakan C\# dalam Unity dengan menerapkan konsep OOP dan design patterns.
  
  \item \textbf{CPMK-4:} Mahasiswa mampu mengimplementasikan sistem input, physics, animation, UI, dan audio dalam game.
\end{enumerate}

\section*{Matriks Keterkaitan CPL dan CPMK}

\begin{table}[h]
\centering
\begin{tabular}{cccc}
  \toprule
  \textbf{CPL} & \textbf{CPMK} & \textbf{Kontribusi} & \textbf{Keterangan} \\
  \midrule
  CPL-1 & CPMK-1 & Tinggi & Penguasaan konsep teoretis game development \\
  CPL-1 & CPMK-2 & Tinggi & Kemampuan menggunakan Unity Engine \\
  CPL-2 & CPMK-3 & Tinggi & Penerapan pemikiran sistematis dalam implementasi C\# \\
  CPL-2 & CPMK-4 & Sedang & Pemikiran kritis dalam implementasi sistem game \\
  CPL-3 & CPMK-2 & Tinggi & Penggunaan tools untuk pengembangan game \\
  CPL-3 & CPMK-3 & Tinggi & Implementasi dengan standar kualitas \\
  CPL-3 & CPMK-4 & Tinggi & Evaluasi kualitas implementasi game \\
  CPL-4 & CPMK-3 & Sedang & Tanggung jawab dalam implementasi kode \\
  CPL-4 & CPMK-4 & Sedang & Kerja sama dalam pengembangan game \\
  \bottomrule
\end{tabular}
\caption{Matriks Keterkaitan CPL dan CPMK}
\end{table}

\cleardoublepage

% ============================================================
% Daftar Isi
% ============================================================
\phantomsection
\addcontentsline{toc}{chapter}{Daftar Isi}
\tableofcontents
\cleardoublepage

\clearpage
\phantomsection
\addcontentsline{toc}{chapter}{Daftar Gambar}
\listoffigures
\cleardoublepage

\phantomsection
\addcontentsline{toc}{chapter}{Daftar Tabel}
\listoftables
\cleardoublepage

\phantomsection
\addcontentsline{toc}{chapter}{Daftar Kode Program}
\lstlistoflistings
\cleardoublepage
